\documentclass[10pt,aspectratio=169]{beamer}

\usepackage[utf8]{inputenc}
\usepackage[T1]{fontenc}
\usepackage[spanish]{babel}
\usepackage{booktabs}
\usepackage{amsmath}
\usepackage{graphicx}
\usepackage{tikz}
\usetikzlibrary{positioning, arrows.meta, shadows, fit, backgrounds}
\usepackage[most]{tcolorbox}
\usepackage{colortbl}

\usetheme{Madrid}
\usecolortheme{default}

\definecolor{uni}{RGB}{122,31,43}
\definecolor{uniink}{RGB}{40,40,40}
\definecolor{job1}{RGB}{31,91,140}
\definecolor{job2}{RGB}{24,124,108}
\definecolor{job3}{RGB}{181,133,36}
\definecolor{etlE}{RGB}{31,91,140}
\definecolor{etlT}{RGB}{24,124,108}
\definecolor{etlL}{RGB}{122,31,43}
\definecolor{dimcolor}{RGB}{24,124,108}
\definecolor{factcolor}{RGB}{122,31,43}

\setbeamercolor{palette primary}{bg=uni,fg=white}
\setbeamercolor{palette secondary}{bg=uni!85!black,fg=white}
\setbeamercolor{palette tertiary}{bg=uni!70,fg=white}
\setbeamercolor{palette quaternary}{bg=uni!40,fg=white}
\setbeamercolor{structure}{fg=uni}
\setbeamercolor{title}{fg=white,bg=uni}
\setbeamercolor{frametitle}{fg=white,bg=uni}
\setbeamercolor{normal text}{fg=uniink}
\setbeamertemplate{navigation symbols}{}
\setbeamersize{text margin left=6mm,text margin right=6mm}

\newcommand{\campos}[1]{\textbf{Campos usados:} #1\par\vspace{2pt}}
\newcommand{\respondea}[1]{\textbf{Responde a:} #1\par\vspace{2pt}}

% caja de dato clave, con color configurable
\newtcolorbox{dato}[2][uni]{
  colback=#1!8, colframe=#1, boxrule=0.8pt, arc=3pt,
  left=6pt, right=6pt, top=4pt, bottom=4pt, title={#2},
  fonttitle=\bfseries\small, coltitle=white, colbacktitle=#1
}

\title[Operaciones policiales -- Spark/PostgreSQL]{Resultados de las operaciones policiales a nivel nacional}
\subtitle{An\'alisis con Apache Spark, Scala y PostgreSQL}
\author[Rojas y Morales]{Luis Antonio Rojas Huaroc \and Breiner Catalino Morales}
\institute[UNI -- CC531 A]{Universidad Nacional de Ingenier\'ia\\ Facultad de Ciencias\\ CC531 A -- An\'alisis en Macrodatos 2026-2}
\date{Laboratorio 03}

\begin{document}

\begin{frame}
\titlepage
\end{frame}

\begin{frame}{Qu\'e exige este laboratorio}
\begin{columns}[T]
\begin{column}{0.5\textwidth}
\begin{dato}[job1]{Dataset}
PNDA, categor\'ia \textbf{Seguridad y Justicia}, no usado por otro grupo.
\end{dato}
\vspace{4pt}
\begin{dato}[job2]{Scala + Spark}
Consultas con \textbf{MapReduce puro} sobre RDD.
\end{dato}
\vspace{4pt}
\begin{dato}[job3]{PostgreSQL}
Arquitectura de tablas propia + proceso \textbf{ETL} documentado.
\end{dato}
\end{column}
\begin{column}{0.5\textwidth}
\begin{dato}[etlL]{Spark SQL}
API de \texttt{DataFrame} (\texttt{groupBy}, \texttt{join}, funciones) + \textbf{vistas temporales}.
\end{dato}
\vspace{4pt}
\begin{dato}[uni]{Excel}
Exportar resultados \textbf{usando Spark}.
\end{dato}
\end{column}
\end{columns}
\end{frame}

\begin{frame}{El dataset: c\'omo lo elegimos}
\begin{dato}[uni]{Primera opci\'on descartada}
\emph{Denuncias policiales registradas a nivel nacional} -- ya estaba tomada por otro grupo.
\end{dato}
\vspace{6pt}
\begin{dato}[job2]{Dataset elegido}
\textbf{Resultados de las operaciones policiales a nivel nacional} [MININTER/CENOPPOL], Observatorio
Nacional de Seguridad Ciudadana. 3 CSV (2023-2025) combinados: \textbf{30\,456 registros}.
\end{dato}
\vspace{6pt}
\begin{dato}[job1]{Columnas}
a\~no, mes, \texttt{VARIABLE} (18 tipos), \texttt{DEPENDENCIAS} (53), \texttt{DIR\_REG\_FRE},
\texttt{VALOR} (decimal real: kilos, soles, personas, unidades).
\end{dato}
\end{frame}

\begin{frame}{Entorno instalado}
\centering
\begin{tikzpicture}[
  caja/.style={draw=job1, fill=job1!10, rounded corners, minimum width=3.6cm, minimum height=1.8cm,
    align=center, font=\small, drop shadow},
  flecha/.style={-{Latex[length=3mm]}, line width=1.2pt, draw=uni}
]
\node[caja] (spark) at (0,0) {\textbf{Apache Spark 3.3.2}\\+ Hadoop 3.3.0\\\textbf{Scala 2.12.15}};
\node[caja, draw=job2, fill=job2!10] (pg) at (5.2,0) {\textbf{PostgreSQL 15.13}\\(binarios port\'atiles)\\+ driver JDBC 42.7.6};
\node[caja, draw=job3, fill=job3!10] (sin) at (10.4,0) {\textbf{Sin privilegios}\\\textbf{de administrador}\\versiones port\'atiles};
\draw[flecha] (spark.east) -- (pg.west);
\draw[flecha] (pg.east) -- (sin.west);
\end{tikzpicture}
\vspace{10pt}

\small Mismas versiones usadas en clase (\texttt{instalacion spark3.3.2.pdf}, \texttt{Semana 11/instalacion.txt}).
\end{frame}

\begin{frame}{Las cinco consultas RDD / MapReduce}
\footnotesize
\begin{tabular}{@{}clp{7.4cm}@{}}
\toprule
\rowcolor{uni} \textcolor{white}{\#} & \textcolor{white}{Tipo} & \textcolor{white}{Pregunta} \\
3.a & MapReduce simple & Mayor/menor valor por tipo de resultado \\
3.b & MapReduce (acumulador) & Promedio, mediana, desviaci\'on, mayor/menor del dinero \\
3.c.1 & \textcolor{job1!80!black}{\textbf{MapReduce triple}} & Desviaci\'on de OPERATIVOS por dependencia vs. promedio nacional \\
3.c.2 & \textcolor{job2!80!black}{\textbf{MapReduce triple}} & \% de cada droga dentro del total de kilos por regi\'on \\
3.c.3 & \textcolor{job3!80!black}{\textbf{MapReduce triple}} & Variaci\'on interanual del dinero incautado promedio mensual \\
\bottomrule
\end{tabular}

\vspace{10pt}
\small \textcolor{job1!80!black}{$\blacksquare$} Job 1: mapea y agrega \quad
\textcolor{job2!80!black}{$\blacksquare$} Job 2: calcula el agregado global \quad
\textcolor{job3!80!black}{$\blacksquare$} Job 3: combina y ordena el resultado final
\end{frame}

\begin{frame}[fragile]{3.a -- Mayor/menor por tipo (MapReduce)}
\respondea{¿Cu\'al es el mayor y menor valor hist\'orico (2023-2025) de cada tipo de resultado?}
\campos{VARIABLE, VALOR}
\begin{verbatim}
val pares = datos.map(_.split(",")).map(c => (c(2), c(5).toDouble))
val maxPorTipo = pares.reduceByKey((a, b) => math.max(a, b))
\end{verbatim}
\begin{dato}[uni]{Resultado}
\footnotesize M\'aximos: DINERO INCAUTADO S/\,4\,028\,271.78; MARIHUANA (KILOS) 29\,924.75\,kg;
CONTRABANDO S/\,513\,820\,390. M\'inimo en los 18 tipos: 0.
\end{dato}
\end{frame}

\begin{frame}[fragile]{3.b -- Estad\'isticos del dinero incautado}
\respondea{¿Promedio, media (mediana), desviaci\'on, mayor y menor del dinero incautado?}
\campos{VARIABLE = "DINERO INCAUTADO (SOLES)", VALOR}
\begin{verbatim}
def combinar(a: Acum, v: Double): Acum =
  (a._1+v, a._2+v*v, a._3+1L, math.max(a._4,v), math.min(a._5,v))
// mediana: sortBy + dos centrales (n par)
\end{verbatim}
\begin{dato}[job2]{Resultado}
\footnotesize $n=1692$ · promedio S/\,31\,840.98 · mediana S/\,1\,799.30 · desviaci\'on
S/\,142\,260.99 · mayor S/\,4\,028\,271.78 (coincide con 3.a) · menor S/\,0.00.
\end{dato}
\end{frame}

\begin{frame}{3.c.1 -- Desviaci\'on de OPERATIVOS por dependencia}
\respondea{¿Qu\'e dependencias reportan muchos m\'as/menos operativos que el promedio nacional?}
\centering
\begin{tikzpicture}[
  job/.style={draw, rounded corners, minimum width=3.5cm, minimum height=1.3cm, align=center,
    font=\scriptsize, drop shadow, text width=3.3cm},
  flecha/.style={-{Latex[length=3mm]}, line width=1.2pt, draw=uniink}
]
\node[job, draw=job1, fill=job1!12] (j1) at (0,0) {\textbf{Job 1}\\suma + conteo\\por dependencia};
\node[job, draw=job2, fill=job2!12] (j2) at (4.6,0) {\textbf{Job 2}\\promedio\\nacional};
\node[job, draw=job3, fill=job3!12] (j3) at (9.2,0) {\textbf{Job 3}\\desviaci\'on\\ordenada};
\draw[flecha] (j1.east) -- (j2.west);
\draw[flecha] (j2.east) -- (j3.west);
\end{tikzpicture}
\vspace{10pt}

\begin{dato}[job3]{Resultado}
\footnotesize Promedio nacional: \textbf{2070.90} operativos/registro. Mayor desviaci\'on:
\textbf{RP Lima} ($+10\,606.94$), RP Cusco ($+6080.13$), RP Ancash ($+3976.21$). DIRFAM: $-2070.90$
(0 operativos).
\end{dato}
\end{frame}

\begin{frame}{3.c.2 -- \% de cada droga por regi\'on/frente}
\respondea{¿Qu\'e droga predomina en cada regi\'on/frente, en kilos?}
\centering
\begin{tikzpicture}[
  job/.style={draw, rounded corners, minimum width=3.5cm, minimum height=1.3cm, align=center,
    font=\scriptsize, drop shadow, text width=3.3cm},
  flecha/.style={-{Latex[length=3mm]}, line width=1.2pt, draw=uniink}
]
\node[job, draw=job1, fill=job1!12] (j1) at (0,0) {\textbf{Job 1}\\suma de kilos\\por (regi\'on, droga)};
\node[job, draw=job2, fill=job2!12] (j2) at (4.6,0) {\textbf{Job 2}\\total de kilos\\por regi\'on};
\node[job, draw=job3, fill=job3!12] (j3) at (9.2,0) {\textbf{Job 3}\\porcentaje\\por regi\'on};
\draw[flecha] (j1.east) -- (j2.west);
\draw[flecha] (j2.east) -- (j3.west);
\end{tikzpicture}
\vspace{10pt}

\small
\begin{tabular}{@{}llrr@{}}
\toprule
\rowcolor{uni} \textcolor{white}{Regi\'on/Frente} & \textcolor{white}{Droga predominante} & \textcolor{white}{Kilos} & \textcolor{white}{\%}\\
DIRNIC & Marihuana & 141\,030.07 & 36.38\%\\
Frente & CC (clorhidrato) & 11\,501.08 & 83.14\%\\
Regi\'on & Marihuana & 91\,960.67 & 82.67\%\\
\bottomrule
\end{tabular}
\end{frame}

\begin{frame}{3.c.3 -- Variaci\'on interanual del dinero incautado}
\respondea{¿Creci\'o o cay\'o el dinero incautado promedio mensual, 2023-2025?}
\centering
\begin{tikzpicture}[
  job/.style={draw, rounded corners, minimum width=3.5cm, minimum height=1.3cm, align=center,
    font=\scriptsize, drop shadow, text width=3.3cm},
  flecha/.style={-{Latex[length=3mm]}, line width=1.2pt, draw=uniink}
]
\node[job, draw=job1, fill=job1!12] (j1) at (0,0) {\textbf{Job 1}\\suma + conteo\\por a\~no};
\node[job, draw=job2, fill=job2!12] (j2) at (4.6,0) {\textbf{Job 2}\\promedio\\mensual};
\node[job, draw=job3, fill=job3!12] (j3) at (9.2,0) {\textbf{Job 3}\\\% variaci\'on\\interanual};
\draw[flecha] (j1.east) -- (j2.west);
\draw[flecha] (j2.east) -- (j3.west);
\end{tikzpicture}
\vspace{10pt}

\small
\begin{tabular}{@{}rrrr@{}}
\toprule
\rowcolor{uni} \textcolor{white}{A\~no} & \textcolor{white}{Suma (S/)} & \textcolor{white}{Prom. mensual (S/)} & \textcolor{white}{\% variaci\'on}\\
2023 & 11\,704\,529 & 975\,377 & --\\
2024 & 25\,880\,606 & 2\,156\,717 & \textcolor{job2!70!black}{\textbf{$+121.12\%$}}\\
2025 & 16\,289\,795 & 1\,357\,483 & \textcolor{uni}{\textbf{$-37.06\%$}}\\
\bottomrule
\end{tabular}
\end{frame}

\begin{frame}{Arquitectura de la base de datos: esquema en estrella}
\centering
\begin{tikzpicture}[
  dim/.style={draw=dimcolor, fill=dimcolor!10, rectangle, rounded corners, minimum width=3.6cm,
    align=left, font=\scriptsize, drop shadow},
  fact/.style={draw=factcolor, fill=factcolor!10, rectangle, rounded corners, minimum width=4cm,
    align=left, font=\scriptsize, drop shadow, line width=1.1pt},
  flecha/.style={-{Latex[length=3mm]}, line width=1.3pt, draw=uniink}
]
\node[dim] (dep) at (0,0) {
  \textbf{\textcolor{dimcolor}{dim\_dependencia}}\\ \underline{dependencia\_id}\\ nombre\\ dir\_reg\_fre};
\node[fact] (hecho) at (5,0) {
  \textbf{\textcolor{factcolor}{hecho\_resultado\_operativo}}\\ \underline{id}\\ anio, mes\\
  dependencia\_id (FK)\\ variable\_id (FK)\\ valor};
\node[dim] (var) at (10,0) {
  \textbf{\textcolor{dimcolor}{dim\_variable}}\\ \underline{variable\_id}\\ nombre\\ categoria\\ unidad};
\draw[flecha] (dep.east) -- (hecho.west) node[midway, above, font=\scriptsize, fill=white] {1:N};
\draw[flecha] (var.west) -- (hecho.east) node[midway, above, font=\scriptsize, fill=white] {1:N};
\end{tikzpicture}
\vspace{10pt}

\small \textcolor{dimcolor}{$\blacksquare$} Dimensiones (53 + 18 filas) \qquad
\textcolor{factcolor}{$\blacksquare$} Tabla de hechos (30\,456 filas)
\end{frame}

\begin{frame}{ETL: Extract -- Transform -- Load}
\centering
\begin{tikzpicture}[
  caja/.style={draw, rounded corners, minimum width=3.6cm, minimum height=2.6cm, align=left,
    font=\scriptsize, text width=3.4cm, drop shadow},
  flecha/.style={-{Latex[length=3mm]}, line width=1.4pt, draw=uniink}
]
\node[caja, draw=etlE, fill=etlE!10] (e) at (0,0) {
  \centering\textbf{\textcolor{etlE}{EXTRACT}}\\[3pt]
  \raggedright 3 CSV de la PNDA (2023-2025), combinados en 30\,456 filas.};
\node[caja, draw=etlT, fill=etlT!10] (t) at (4.6,0) {
  \centering\textbf{\textcolor{etlT}{TRANSFORM}}\\[3pt]
  \raggedright Correcci\'on de 432 filas; \texttt{categoria}/\texttt{unidad} derivadas; llaves con
  \texttt{row\_number()}.};
\node[caja, draw=etlL, fill=etlL!10] (l) at (9.2,0) {
  \centering\textbf{\textcolor{etlL}{LOAD}}\\[3pt]
  \raggedright 3 tablas escritas a PostgreSQL por JDBC (\texttt{df.write.format("jdbc")}).};
\draw[flecha] (e.east) -- (t.west);
\draw[flecha] (t.east) -- (l.west);
\end{tikzpicture}
\end{frame}

\begin{frame}{Evidencia: ETL y base de datos funcionando}
\centering
\includegraphics[width=0.78\linewidth]{db_tablas_conteo.png}
\end{frame}

\begin{frame}{Spark SQL -- API de DataFrame (5.a -- 5.b)}
\begin{dato}[job1]{5.a -- groupBy + filter (4 consultas)}
Suma por (a\~no,categor\'ia); conteo por (dependencia,mes); kilos por (regi\'on,categor\'ia); total por
(a\~no,unidad).
\end{dato}
\vspace{10pt}
\begin{dato}[job2]{5.b -- promedio + orderBy}
Promedio de \texttt{valor} por categor\'ia; dependencias ordenadas por suma total desc (DIRPOFIS
S/\,629.8M lidera).
\end{dato}
\end{frame}

\begin{frame}{Spark SQL -- API de DataFrame (5.c -- 5.d)}
\begin{dato}[job3]{5.c -- los 4 tipos de join}
\texttt{inner}/\texttt{left}/\texttt{right}/\texttt{full} -- los 4 devuelven 30\,456 filas,
\textbf{0 registros hu\'erfanos}.
\end{dato}
\vspace{10pt}
\begin{dato}[uni]{5.d -- functions.\_ (4 consultas)}
\texttt{round(avg)}; \texttt{when/otherwise} en rangos; \texttt{concat\_ws}; \texttt{stddev} por
categor\'ia.
\end{dato}
\end{frame}

\begin{frame}{Vistas temporales + SQL puro (5.e -- 5.f)}
\begin{dato}[job1]{5.e -- 3 joins (objetivos distintos a 5.c)}
Dinero por dependencia+a\~no (INNER); operativos por dependencia (LEFT); variables sin uso en 2025
(FULL -- ninguna).
\end{dato}
\vspace{10pt}
\begin{dato}[job2]{5.f -- 3 ORDER BY + filtro}
Armas de mayor a menor; regi\'on 2025 de mayor a menor; dinero $>0$ de menor a mayor.
\end{dato}
\end{frame}

\begin{frame}{Exportaci\'on a Excel con Spark (5.g)}
\begin{dato}[job3]{5.g -- Excel con Spark}
Resultados de 5.e y 5.f exportados con \texttt{spark-excel}, una hoja por consulta, 2 archivos de 3
hojas cada uno.
\end{dato}
\vspace{10pt}
\centering
\includegraphics[width=0.75\linewidth]{img_excel.png}
\end{frame}

\begin{frame}{Conclusiones (1/2)}
\begin{dato}[uni]{1}
5 consultas RDD/MapReduce + 18 consultas Spark SQL (DataFrame + vistas temporales), todas sobre 30\,456
registros reales de operaciones policiales 2023-2025.
\end{dato}
\vspace{10pt}
\begin{dato}[job1]{2}
Esquema en estrella en PostgreSQL poblado por un ETL completo en Scala/Spark, con correcci\'on de un
problema real de calidad de datos.
\end{dato}
\end{frame}

\begin{frame}{Conclusiones (2/2)}
\begin{dato}[job2]{3}
Verificaci\'on cruzada: el m\'aximo de dinero incautado coincide exacto entre dos consultas
independientes; los 4 joins no encuentran registros hu\'erfanos.
\end{dato}
\vspace{10pt}
\begin{dato}[job3]{4}
El dinero incautado promedio mensual creci\'o 121\,\% en 2024 y cay\'o 37\,\% en 2025 -- alta
volatilidad, consistente con su desviaci\'on est\'andar.
\end{dato}
\end{frame}

\end{document}
